\booksection{제11장 부록:전문가용 이미징기술 지식}

\chapterintro{스마트폰 사진과 영상의 진정한 오묘함을 이해하려면, 일부 전문 영상 기술 지식을 알아야 합니다. 이러한 지식은 스마트폰 카메라를 더 잘 활용할 뿐만 아니라, 기기 선택 시 판단력을 키워줍니다.하나, 해상도, 프레임 레이트와 비트레이트: 영상 품질의 3대 요소해상도(Resolution):해상도는 화면의 선명도를 결정하며, 픽셀 수로 표시됩니다.\textbf{720P (HD)}: 1280×720 픽셀, 입문급 화질.\textbf{1080P (Full HD)}: 1920×1080 픽셀, 현재 가장 주류인 해상도.\textbf{2K (QHD)}: 2560×1440 픽셀, 1080P보다 더 선명함.\textbf{4K (UHD)}: 3840×2160 픽셀, 초고화질, 디테일이 풍부함.\textbf{8K}: 7680×4320 픽셀, 미래 표준, 현재 기기 지원은 제한적.해상도가 높을수록 화면이 선명해지지만, 파일도 커지고 저장 및 처리 능력 요구사항이 높아집니다.프레임률(Frame Rate):프레임률은 화면의 부드러움을 결정하며, 초당 프레임 수(fps)로 표시합니다.\textbf{24fps}: 영화 표준 프레임률, 독특한 “영화감”이 있습니다.\textbf{30fps}: TV와 네트워크 영상에서 자주 사용, 부드럽고 자연스러움.\textbf{60fps}: 고 프레임률, 운동 장면과 게임 녹화에 적합, 매우 부드러움.\textbf{120fps/240fps}: 초고 프레임률, 슬로우 모션 영상에 사용 가능.프레임 레이트가 높을수록 화면이 더 부드럽고 슬로우 모션 효과가 좋지만, 파일도 커집니다.비트레이트(Bitrate):비트레이트는 영상의 데이터 밀도를 결정하며, 초당 데이터량(Mbps)으로 표시됩니다.비트레이트가 높을수록 화질이 좋고 디테일이 더 보존되지만 파일도 커집니다.비트레이트가 너무 낮으면 화면에 모자이크, 색밴딩 등 압축 흔적이 나타납니다.전문가 영상 촬영은 일반적으로 높은 비트레이트가 필요하여 후반 처리 여지를 더 보존합니다.이,영상 코덱 형식:H.264, H.265와 ProResH.264 (AVC):\textbf{특징}: 가장 널리 사용되는 코덱 형식, 호환성이 매우 좋음.\textbf{장점}: 거의 모든 기기에서 재생 가능, 압축 효율은 중간 정도.\textbf{단점}: 새 형식에 비해 압축 효율이 낮고 파일이 큼.\textbf{적용 시나리오}: 네트워크 영상 공유, 일반 사용자 일상 촬영.H.265 (HEVC):\textbf{특징}: H.264의 업그레이드 버전, 압축 효율 약 50\% 향상.\textbf{장점}: 같은 화질에서 파일이 더 작고, 더 높은 해상도와 프레임률 지원.\textbf{단점}: 코덱과 디코딩에 더 강력한 프로세서가 필요하며, 일부 구형 기기는 호환되지 않음.\textbf{적용 시나리오}: 4K 영상 촬영, 저장 공간이 제한된 상황.Apple ProRes:\textbf{특징}: 애플이 개발한 전문가 영상 코덱, 거의 무손실 압축.\textbf{장점}: 화질이 매우 높고, 후반 처리 여지가 크며, 색이 정확함.\textbf{단점}: 파일이 매우 크고, 대량의 저장 공간이 필요하며, 처리에 전문가 기기가 필요함.\textbf{적용 시나리오}: 전문가 영상 제작, 영화급 촬영.ProRes RAW:\textbf{특징}: ProRes와 RAW 형식의 장점을 결합.\textbf{RAW란}: 처리되지 않은 원본 센서 데이터, 최대 정보량을 보존.\textbf{장점}: 후반 조정 여지가 매우 크고, 노출, 화이트 밸런스 등 매개변수를 화질 손실 없이 변경 가능.\textbf{단점}: 파일이 매우 커서 전문가 후반 소프트웨어로 처리해야 함.\textbf{적용 시나리오}: 고급 전문가 제작. iPhone 15 Pro 시리즈 이후에서 Apple ProRes 영상 촬영 지원(최대 4K 60fps), iPhone 17 Pro는 한 단계 더 나아가 ProRes RAW 기기 내 직접 녹화 지원.삼,HDR: 고 다이내믹 레인지 기술HDR(High Dynamic Range) 기술은 화면의 가장 밝은 부분과 가장 어두운 부분의 디테일을 동시에 보존할 수 있습니다.\textbf{HDR10}: 가장 기초적인 HDR 표준, 호환성이 좋음.\textbf{HDR10+}: 삼성이 추진한 표준, 동적 메타데이터 지원.\textbf{Dolby Vision}: 돌비가 개발한 최고급 HDR 표준, 화질이 가장 좋지만 라이선스가 필요함.\textbf{HLG}: 방송급 HDR 표준, 일반 디스플레이와 호환.HDR 영상을 촬영할 때는 HDR을 지원하는 디스플레이가 있어야 완전한 효과를 볼 수 있습니다.사,색 공간과 색 심도색 공간:\textbf{sRGB}: 표준 색 공간, 호환성이 가장 좋음.\textbf{Display P3}: 애플이 보급한 광색역 표준, 색이 더 풍부함.\textbf{Rec.2020}: HDR 영상의 표준 색 공간, 색역이 가장 넓음.색심도:\textbf{8-bit}: 1670만 색, 일반 영상 표준.\textbf{10-bit}: 10.7억 색, 전문가 영상 표준, 색 전환이 더 부드러움.\textbf{12-bit}: 680억 색, 고급 전문가 표준.실용 권장일상 촬영: 1080P/4K + 30fps + H.265 코덱, 화질과 파일 크기의 균형.운동 촬영: 1080P/4K + 60fps, 부드러움 보장.전문가 창작: 4K + ProRes(지원하는 경우), 최대 후반 여지 보존.저장 권장: 높은 비트레이트 영상은 공간을 많이 차지하므로 정기적으로 컴퓨터나 클라우드 디스크에 백업할 것을 권장.한번 탐색해 보세요: 당신의 스마트폰 영상 설정카메라 App을 열고, 영상 모드로 전환하세요.해상도와 프레임률 옵션을 확인하고, 다양한 조합을 시도해 보세요.지원하는 경우, HDR, ProRes 등 고급 옵션이 있는지 확인하세요.다양한 설정으로 영상을 몇 장 촬영하고, 화질과 파일 크기를 비교하세요.오,RAW 형식과 HEIF/JPEG 비교RAW 형식이란:RAW 형식은 카메라 센서가 기록한 “원본 데이터”입니다—어떠한 압축과 처리도 거치지 않은 순수 원본 정보입니다. 신선한 식재료와 같아서 후반 “요리”가 필요합니다. 스마트폰은 저장과 연산 능력이 제한되어 기본값으로 촬영된 것은 이미 “요리된” JPEG 또는 HEIF이며, RAW는 별도로 켜야 합니다.RAW vs HEIF vs JPEG 비교:\textbf{JPEG}: 가장 오래된 형식, 호환성이 가장 좋음(어떤 기기에서도 열림). 8-bit 색 심도, 손실 압축, 파일이 가장 작음(사진 한 장 2-5MB). 후반 여지가 제한적이며, 어두운 부분을 억지로 밝게 하면 색밴딩이 나타나기 쉽습니다.\textbf{HEIF/HEIC}: 애플이 주도하는 새 형식, 같은 화질에서 파일 크기가 JPEG의 절반(사진 한 장 1-3MB). 10-bit 색 심도와 투명 채널 지원. 단점은 일부 구형 기기와 Windows는 플러그인을 설치해야 열 수 있습니다.\textbf{RAW}: 압축되지 않은 원본 데이터, 파일이 가장 큼(사진 한 장 20-100MB). 색 심도는 일반적으로 12-16bit, 후반 여지가 가장 큼. 단점은 전문가 소프트웨어(Lightroom, Snapseed 등)로만 열고 처리할 수 있으며, 직접 공유할 수 없습니다.전문가 사진사가 RAW를 찍는 이유:\textbf{화이트 밸런스를 후반에 무손실 조정 가능}: JPEG에서 화이트 밸런스를 잘못 맞추면(예: 백열등 아래에서 누렇게), 후반 보정은 화질 손실을 초래합니다; RAW는 색온도를 다시 설정하여 완전히 무손실로 보정할 수 있습니다.\textbf{노출 허용 오차가 큼}: JPEG에서 2스톱 노출 부족인 어두운 부분은 거의 새까맣고, 밝게 하면 노이즈투성이; RAW는 3-4스톱 노출 부족인 어두운 부분에서도 사용 가능한 디테일을 “구출”할 수 있습니다. 과노출도 마찬가지.\textbf{어두운 부분과 밝은 부분의 디테일이 더 풍부함}: RAW의 다이내믹 레인지는 JPEG보다 훨씬 넓어, 밝은 부분이 과노출되지 않고 어두운 부분이 새까맣지 않은 여지가 훨씬 큽니다.\textbf{선명화, 노이즈 감소, 렌즈 보정 모두 후반에 제어 가능}: 카메라의 기기 내 처리는 “일방적”이지만, RAW는 반복적으로 조정할 수 있고, 매번 원본 데이터에서 시작합니다.스마트폰에서 RAW 켜는 법:\textbf{iPhone ProRAW}: iPhone 12 Pro 이상 Pro 모델만 지원. “설정 카메라 형식”에서 “Apple ProRAW”를 켜세요. 그러면 카메라 인터페이스 상단에 “RAW” 버튼이 나타나며, 켜기만 하면 됩니다. 파일은 약 25MB/장이며, 필요할 때만 켜는 것을 권장합니다.\textbf{안드로이드 Motion Raw / Ultra RAW}: 브랜드마다 부르는 이름이 다름. 샤오미/OPPO/vivo 플래그십 모델은 카메라 “전문가 모드” 또는 “더 보기”에서 RAW 켜기/끄기를 찾을 수 있습니다; 화웨이는 “전문가 모드”에서 RAW를 켭니다; 삼성 Galaxy S 시리즈는 “전문가” 모드에 RAW 옵션이 있습니다. 파일 크기는 일반적으로 20-50MB.\textbf{일반 권장}: 일상 촬영에서는 RAW를 켤 필요가 없으며, 파일이 너무 커서 공간을 차지합니다; 중요한 장면(풍경, 인물, 후반 필요)에서만 켜세요.육,컴퓨테이셔널 포토그래피 원리현대 스마트폰 사진은 “컴퓨테이셔널 포토그래피”없이 할 수 없습니다—즉, 렌즈와 센서뿐만 아니라 알고리즘에 의존해 사진을 생성합니다. 이 몇 가지 핵심 원리를 이해하면, 셔터를 누르는 순간 스마트폰이 무엇을 하는지 더 잘 알 수 있습니다.다중 프레임 합성:스마트폰에서 셔터를 한 번 누르면, 실제로는 여러 장의 사진을 찍습니다. 알고리즘이 다중 프레임을 정렬하고, 가중 융합하여 한 장의 고품질 결과 이미지를 얻습니다.\textbf{야간 모드}: 1-3초 내에 여러 장의 노출 부족 사진을 연속 촬영하고, 정렬 후 합성—노이즈는 서로 상쇄되고, 밝기는 누적됩니다. 결과적으로 어두운 환경에서도 깨끗한 화면을 촬영할 수 있습니다.\textbf{HDR(고 다이내믹 레인지)}: 노출 부족, 정상, 과노출 세 장(또는 그 이상)을 각각 촬영하여, 밝은 부분과 어두운 부분 모두 디테일이 있는 사진 한 장으로 합성. 역광 장면에 특히 유효.\textbf{Deep Fusion(애플)/Photonic Engine(광자 엔진)}: 중저조도 장면에 대해, 픽셀 단위로 다중 프레임을 분석하고 융합하며, 텍스처와 디테일을 중점적으로 향상시켜, 옷, 머리카락 같은 디테일이 눈에 띄게 더 선명해집니다.의미 분할:AI는 단순히 전체 이미지를 일괄 처리하는 것이 아니라, 먼저 화면에 무엇이 있는지 “이해”합니다—하늘, 얼굴, 잔디, 건물, 옷—그런 다음 각 영역에 다른 처리 전략을 적용합니다.\textbf{하늘}: 파란색을 개별적으로 강화하고, 밝기를 낮추어 과노출 방지.\textbf{얼굴}: 피부색 영역을 인식하고, 개별적으로 밝게 하고, 피부를 부드럽게 하며, 오관과 머리카락 디테일은 건드리지 않음.\textbf{잔디/나뭇잎}: 녹색 채도를 강화하고, 디테일 선명화를 향상.\textbf{배경 흐림}: 인물 모드에서 AI가 피사체 윤곽을 인식하고, 배경은 흐리게, 피사체는 선명하게.딥러닝의 활용:\textbf{AI 장면 인식}: 카메라가 자동으로 “음식”“반려동물”“야경”“문서” 등 장면을 인식하고, 해당 최적화 매개변수를 적용. 예를 들어 음식 모드는 웜톤과 채도를 강화하여 음식이 더 먹음직스럽게 보이게 합니다.\textbf{AI 인물 흐림}: 스마트폰 렌즈는 물리적으로 DSLR 같은 얕은 피사계 심도를 구현할 수 없어, 딥러닝 모델로 피사체 가장자리를 예측하고, 배경을 흐리게 합니다. 최신 알고리즘은 머리카락 가닥, 유리 가장자리 같은 디테일도 올바르게 처리할 수 있습니다.\textbf{AI 제거/AI 확장}: 화면의 행인이나 잡물을 선택하면, AI가 문맥을 기반으로 가려진 부분을 “보완”합니다. 이는 생성형 AI가 이미지 분야에 활용된 것입니다.컴퓨테이셔널 포토그래피 vs 광학 포토그래피:\textbf{컴퓨테이셔널 포토그래피의 장점}: 작은 크기로도 좋은 화질을 얻을 수 있고, 야간, 역광 등 극한 장면에서도 촬영 가능하며, 일반 사용자가 사진 지식 없이도 “멋진” 사진을 찍을 수 있습니다.\textbf{컴퓨테이셔널 포토그래피의 단점}: 화면이 때때로 “과도하게 처리”됨—과도한 선명화, 색 왜곡, 부자연스러운 하늘 색; 다중 프레임 합성은 피사체 이동 시 잔상이 남음; 디테일은 “계산된” 것이지 실제 기록이 아님.\textbf{광학 사진의 장점}: 큰 센서+큰 조리개 렌즈가 빛을 진짜로 포착하여, 화질이 자연스럽고 후반 여지가 큼; 알고리즘 “추측”으로 인한 아티팩트가 발생하지 않음.\textbf{트렌드}: 현대 플래그십 스마트폰은 “광학+컴퓨테이션” 결합—큰 센서로 실제 정보를 획득하고, 알고리즘으로 최적화하여 표시. 둘은 대립하지 않습니다.칠,렌즈 유형과 초점 거리현대 플래그십 스마트폰은 일반적으로 3-4개 렌즈가 있으며, 각각 다른 초점 거리에 해당합니다. 초점 거리를 이해하면, 다른 장면에서 올바른 렌즈를 선택할 수 있습니다.일반적인 렌즈와 등가 초점 거리:\textbf{초광각(13-24mm 등가)}: 시야가 가장 넓고, 풍경, 건물, 실내 공간에 적합. 눈앞의 대장면을 촬영할 수 있지만, 가장자리에 왜곡이 있습니다(직선이 휘어짐). 인물 촬영 시 얼굴이 늘어납니다. “대장면”에 적합, “클로즈업”에는 부적합.\textbf{주 카메라(24-28mm 등가)}: 일상에서 가장 자주 사용하는 렌즈, 화질이 가장 좋고, 조리개가 가장 큼. 시야와 화질을 모두 고려하며, 인물과 풍경 모두 촬영 가능. 대부분의 스마트폰 “1x”가 주 카메라입니다.\textbf{망원(70-200mm 등가)}: 원경, 인물 클로즈업 촬영에 사용. 망원은 공간감을 압축할 수 있어(멀리 있는 산과 가까이 있는 사람이 더 가까이 보임), 인물 촬영 시 더 자연스러운 오관 비율과 배경 흐림을 얻을 수 있습니다.\textbf{잠망경식 망원(200mm+ 등가)}: 반사식 광로 설계로, 스마트폰 두께 내에서 더 긴 망원 거리를 구현. 아주 먼 거리에서도 디테일을 선명하게 촬영할 수 있어, 콘서트, 스포츠 이벤트, 야생 동물에 적합. 플래그십 모델에만 있습니다.센서 크기와 픽셀 크기의 트레이드오프:센서는 “베이스”이며, 받을 수 있는 빛의 양을 결정합니다. 업계에서 “큰 베이스가 한 급 차이 나면 압도한다”고 합니다—같은 픽셀 수에서, 센서가 클수록 단일 픽셀도 커지고, 빛 수용량이 많아지며, 화질이 좋아집니다.\textbf{큰 센서}: 빛 수용량이 많고, 야경 화질이 좋으며, 다이내믹 레인지가 넓음. 단점은 스마트폰이 두꺼워야 하고(렌즈 모듈 돌출), 비용이 높음.\textbf{고화소 작은 센서}: 낮에 빛이 충분할 때, 고화소는 더 선명한 사진을 촬영할 수 있어 후반 크롭에 편리. 하지만 야간에는 각 픽셀의 빛 수용량이 부족하여 노이즈가 많음.\textbf{주류 플래그십 선택}: 주 카메라는 대부분 1/1.3인치에서 1인치(일부 모델은 1인치 대형 센서 사용), 화소 수 5000만 정도—이것이 화질과 크기의 균형점입니다.픽셀 병합 기술(Pixel Binning):낮의 고화소와 야간의 대화소를 모두 고려하기 위해, 스마트폰 제조사는 “픽셀 병합”을 사용합니다—여러 작은 픽셀을 하나의 큰 픽셀로 합쳐서 출력.\textbf{사병합 픽셀(Quad Bayer)}: 5000만 화소 센서, 1250만 화소 사진 출력. 각 출력 픽셀은 4개의 인접 픽셀로 합성되어, 빛 수용량이 4배 향상되고, 화질이 더 깨끗함.\textbf{구병합 픽셀(Nonacell)}: 1억 화소 이상 센서에 사용, 약 1100-1200만 화소 출력. 같은 원리로, 빛 수용량이 더 큼.\textbf{실제 효과}: 기본 촬영 모드에서, 당신이 받는 것은 병합 후 저화소 사진입니다—화질이 좋고, 파일이 작음. 극한 디테일이 필요할 때(예: 문서, 원경 촬영), 수동으로 “고화소 모드”로 전환하여 전체 해상도로 출력할 수 있습니다.팔,특수 촬영 모드일상 촬영 외에도, 스마트폰에는 몇 가지 특수 모드가 있으며, 올바르게 사용하면 평소에는 촬영할 수 없는 화면을 담을 수 있습니다.파노라마 모드:촬영 시 화살표를 따라 스마트폰을 천천히 안정적으로 가로로 이동하고, 빠르게 하지 마세요—빠르면 이어 붙이기 어긋남이 발생합니다.스마트폰을 수직 방향으로 안정적으로 유지하고, 가로로만 평행 이동. 상하로 흔들리면 지평선이 물결모양이 됩니다.이동하는 사람이나 차량에는 신중하게 사용—이어 붙일 때 “반쪽 몸”이나 “유령”이 나타날 수 있습니다.파노라마 모드는 세로로도 촬영할 수 있습니다—세로 촬영 가로 이어 붙이기, 높이 솟은 건물이나 폭포 촬영에 적합.슬로우 모션과 타임랩스:\textbf{슬로우 모션}: 120fps 또는 240fps 고 프레임률로 촬영하고, 재생 시 슬로우 모션이 됨. 물방울, 운동 순간, 반려동물 점프 촬영에 적합. 빛이 충분해야 하며, 슬로우 모션 모드에서는 빛 수용 시간이 짧아 야간 화질이 떨어짐.\textbf{타임랩스}: 간격을 두고 여러 장의 사진을 촬영해 빨리 감기 영상으로 합성. 차량 흐름, 구름 이동, 일몰, 꽃 피는 모습 촬영에 적합. 촬영 시간은 몇 분에서 몇 시간까지이며, 반드시 삼각대로 스마트폰을 고정해야 하며, 손에 들면 반드시 흐려짐.\textbf{소팁}: 타임랩스로 구름 촬영 시, 스마트폰을 설치하고 15-30분 촬영하면, 마지막에 10초 정도의 타임랩스 영상을 얻을 수 있으며, 효과가 압도적입니다.별하늘/천문 모드:도시 빛 공해에서 멀리 떨어져, 완전히 가로등 간섭이 없는 교외나 산악 지대를 찾아야. 도시에서는 은하를 촬영할 수 없습니다.반드시 삼각대를 사용하고, 손에 들면 반드시 흐려짐. 스마트폰을 몇 초에서 수십 초 동안 안정적으로 설치해야 함.날씨는 맑고 구름이 없어야 하며, 달이 너무 크지 않아야(보름달은 별빛을 가림). 음력 월초나 월말의 어두운 밤이 가장 적합.iPhone의 “나이트 모드”는 자동으로 별하늘을 인식하고 노출 시간을 연장(최대 30초); 안드로이드 플래그십의 “별하늘 모드” 또는 “천문 모드”도 비슷한 효과. 인내심을 가지고, 셔터를 누른 후 스마트폰을 움직이지 마세요.접사 모드:꽃, 곤충, 보석, 문자 클로즈업 등 작은 물체 촬영에 사용. 1-2cm 거리까지 초점을 맞출 수 있음.일부 스마트폰에는 전용 접사 렌즈가 있음(대부분 200만 또는 500만 화소의 소형 센서 렌즈)으로 화질이 보통; 플래그십 모델은 초광각 렌즈로 접사를 겸하여, 화질이 더 좋음.접사 촬영 시 렌즈가 피사체에 매우 가까워 빛을 가릴 수 있으며, 필요시 탁자 조명이나 손전등으로 보조 조명.초점 거리가 매우 가까워 피사계 심도가 매우 얕고, 피사체 전후로 몇 밀리미터만 이동해도 초점이 벗어날 수 있음. 촬영 시 손을 안정적으로 유지하고, 피사체를 움직이지 않게 하는 것이 좋습니다.정적 사진 실용 권장일상 촬영: HEIF/JPEG 기본값이면 충분, 파일이 작고 화질이 충분함.중요 장면(풍경, 인물, 후반 필요): RAW를 켜서 최대 후반 여지를 보존.렌즈 선택: 풍경은 초광각, 인물은 망원, 일상은 주 카메라.야경과 별하늘: 반드시 스마트폰을 고정하고, 야간 모드 또는 천문 모드를 켜세요.대형 센서 스마트폰(예: 1인치 센서 모델)은 화질 상한이 더 높지만, 조작 습관이 여전히 관건.직접 해보기: 당신의 카메라 능력 파악하기스마트폰이 RAW 촬영을 지원하는지 확인(설정카메라형식), 켜고 한 장 촬영해 보고, 같은 장면의 JPEG과 RAW 파일 크기를 비교하세요.초광각, 주 카메라, 망원 세 가지 초점 거리로 전환하여, 같은 원경을 각각 촬영하고, 시야 차이와 화질 차이를 체험하세요.야간 모드 또는 별하늘 모드가 있는 장면을 찾아, 스마트폰을 고정한 후 야경 사진을 한 장 촬영하세요.파노라마 모드로 가로 촬영을 한 장 해보고, 스마트폰을 안정적으로 이동하는 데 주의하세요.}
